\section*{Part VIII: A Founder's Method for Producing Meaningful Capability Improvement}

\subsection*{112. Begin with a Consequential Performance, Not a Content Inventory}

Write down 1 outcome that a real learner or sponsoring organization would value. The sentence must name the \textbf{actor, task family, context, performance standard, and permitted resources}. ``Understand machine learning'' is too vague. ``Given an unfamiliar predictive workflow, identify leakage and defend an evaluation design under realistic deployment constraints'' is testable.

Then ask what happens if the learner never acquires the capability. Is the loss significant enough for someone to protect time or pay? This is the business question. Separately ask whether the competence is learnable through the resources you can responsibly offer. That is the educational question.

Do not assume that an important human problem automatically produces a viable education product. Discovery and learning design must inform each other without collapsing into 1 discipline.

\subsection*{113. Diagnose the Learner's Current Performance Through Authentic Tasks}

Before building the course, obtain samples of work from representative members of the intended entry population, with informed consent and appropriate privacy safeguards. Give an unfamiliar task of manageable scope. Ask participants to explain their reasoning, not just provide the final answer.

Record the distribution of error \emph{types}: missing prerequisite knowledge, incorrect mental model, inability to choose a method, fragile execution, weak verification, or poor communication. A useful diagnosis identifies not merely who scored low but \textbf{why they could not perform}.

If the task is too hard or too easy for nearly everyone, fix the diagnostic before drawing conclusions about the market. A test that cannot distinguish learner states is a poor instrument for educational design.

\subsection*{114. Construct the Dependency Graph of the Capability}

Starting from the target performance, identify which concepts, procedures, decisions, and perceptual cues are required. Separate \textbf{essential prerequisites}, \textbf{supporting knowledge}, and \textbf{interesting extensions}. Look for dependencies where a small misunderstanding causes many later errors.

A founder teaching applied AI may find that evaluation design is a bottleneck for many tasks. Students who cannot distinguish valid and invalid data splits cannot safely interpret the performance of a more sophisticated model. That bottleneck deserves repeated practice and scrutiny.

The point is not to create the largest prerequisite graph. It is to find a teachable path through the smallest number of consequential dependencies.

\subsection*{115. Choose the Learning Mechanism That Matches the Failure}

Different failure types require different interventions:

\begin{table*}[t]
\centering
\small
\caption{Matching the Learning Mechanism to the Diagnosed Failure Type}
\begin{tabularx}{\textwidth}{p{0.25\textwidth}p{0.38\textwidth}X}
\toprule
\textbf{Observed Failure} & \textbf{Likely Next Learning Activity} & \textbf{Why It May Help} \\
\midrule
Cannot recall a necessary fact & Brief explanation and spaced retrieval & Establishes and makes knowledge available \\
Misunderstands a concept & Contrasting examples and causal explanation & Repairs the mental model \\
Can follow but not construct & Worked example with gradually removed steps & Builds independent procedure \\
Uses wrong method & Mixed and unlabeled cases & Trains conditional selection \\
Fails on messy cases & Increasingly authentic scenarios & Builds robust application \\
Repeats a subtle error & Diagnostic feedback followed by retry & Repairs the specific weakness \\
Avoids work or cannot sustain it & Time, relevance, and support diagnosis & Addresses participation conditions \\
Cannot transfer & Changed-context tasks and explicit abstraction & Trains retrieval in new settings \\
Cannot judge AI output & Verification exercises and adversarial cases & Builds evidence-sensitive oversight \\
\bottomrule
\end{tabularx}
\end{table*}

These are \textbf{design heuristics} supported by varying degrees of relevant research, not guarantees. The founder should observe whether the chosen intervention changes the next independent attempt.

\subsection*{116. Write the Attempt-Feedback-Retry Loop Explicitly}

Every core skill should have at least 1 place where the learner attempts work before receiving the complete solution, receives meaningful information about the attempt, and tries again in a way that tests improvement.

The task might be a 2-minute calculation, a 30-minute debugging problem, a design memo, or a multi-day experiment. The appropriate scale depends on the competence. It should be possible for the instructor to say exactly \textbf{what was attempted, what the result revealed, and what the learner changed}.

An assignment with no feedback may be practice but cannot reliably support targeted correction. Feedback with no opportunity to revise is a report, not a complete learning cycle. A correction made only by the instructor is not evidence that the learner can now perform it.

\subsection*{117. Design Delayed and Changed-Context Reuse from the Beginning}

A lesson should specify where its essential knowledge reappears. Some ideas need retrieval days later; others need integration into subsequent modules; some need a more difficult transfer task at program end.

If a concept is crucial for a long-term profession, there should be an opportunity to use it after the original teaching context is gone. This does not require assessing every fact at every interval. Select a small number of essential capabilities and design representative delayed prompts.

The easiest time to make space for retention is \textbf{before the syllabus becomes overcrowded}. Once every week is committed to a new technology, retrieval and practice are treated as optional extras. That is a structural error, not a learner's character flaw.

\subsection*{118. Make the Level of Independence Explicit}

For each task, specify whether learners may consult notes, documentation, peers, AI tools, templates, or instructors. Some tasks should be genuinely independent. Others should reflect the collaborative and tool-rich environment of work.

Do not use ``independent'' as an emotional adjective. State the actual restrictions and the judgment required. An engineer who uses documentation intelligently is not necessarily less competent than 1 who memorizes an API; an engineer whose AI wrote a system they cannot verify is not necessarily competent simply because the system runs.

2 conditions often strengthen inference: a controlled check of foundational reasoning and an authentic assisted performance with clear responsibility for correctness.
